% Page 014

\seite{14}

\margmark{Dienstwohnung!}%
war früher Wohnung des Vikars Peter Gillesgord\unsicher,\ob
welcher nach Amerika auswanderte. Die Gemeinde\ob
kaufte es im Jahre 1843, ließ abändern und\ob
später noch nötige Vergrößerungsarbeiten ausführen,\ob
trotzdem entspricht es den heutigen Anforderungen\ob
nicht mehr. Und der Plan, ein neues Schulhaus\ob
zu bauen, ist bereits seit einige Jahre aufge\ohb
nommen worden, ohne jetzt zur Verwirk\ohb
lichung zu kommen. Es mangelt nämlich\ob
nur an dem nötigen Baugeld.\ob
\unsicher ist entstanden, und das\unsicher.

\pend
\begin{verse}
Im Jahre 1886.\ war es!
\end{verse}
\pstart

Die Schülerzahl war im hiesigen Jahre 44.\ Die fand sich\ob
gänzlich befriedigend nicht. In diesem Jahre wurde eine\ob
kleine Reparatur an dem Schulhause vorgenommen.\ob
Da kürzlich befand sich früher im Formweg Thüren\ob
wurde besetzt indem die Schule gelegt und großer\ob
gebrochen nun eine Wörte des Garthaus, welcher\ob
\margmark{Weltlicher Gebrauch\\beim\\Lehrer Brinn [?].}%
dem Lehrer hier à 3 M.\ angewiesen ist,\ob
hatte der Winterpracht hereingetragen. Im Oder\unsicher\ob
Genuß der üblichen Gebräuche, dem Mittlinger\unsicher, der\ob
sich ein Theil beschafft, sämtliche Fuhren\ob
für Wein, Brot, Holz und Licht Bedürfniß zu decken.\ob
Die Ländereien werden meistens von Ortseinwohnern bebaut\ob
und zu diesem Zwecke werden dem \enquote{Leutenden} auch die\ob
nötigen Schulmittel gebracht. Dieser edle Zug\ob
erinnert an die alten kirchlichen Einrichtungen.
\pend
